\chapter{Decision Trees}
\companion{StatQuest: Decision and Classification Trees, Clearly Explained}{\ytid{_L39rN6gz7Y}}

Decision trees are the building block of random forests, gradient boosting, XGBoost and LightGBM, which together win
most tabular problems, including most at InfoEdge. Glassdoor reports for InfoEdge include ``pseudo-code a decision
tree''; the experience notes include reducing overfitting in trees. Be able to compute impurity by hand and write the
algorithm from memory.

\section{A tree is a game of twenty questions}
To guess whether a seeker will apply to a job, ask a sequence of yes/no questions:
\begin{center}
\begin{tikzpicture}[level distance=1.25cm,sibling distance=4.2cm,
  every node/.style={draw,rounded corners,font=\small,align=center,fill=accentlight},
  edge from parent/.style={draw,-{Stealth}}]
\node {Job in preferred city?}
  child {node {Salary $\ge$ expected?}
    child {node[fill=exbg] {Apply: 72\%}}
    child {node[fill=trbg] {Apply: 18\%}}}
  child {node {Remote?}
    child {node[fill=exbg] {Apply: 40\%}}
    child {node[fill=trbg] {Apply: 4\%}}};
\end{tikzpicture}
\end{center}
\begin{itemize}
\item Each internal \textbf{node} tests one feature against a threshold (``experience $\le 3.5$'') or a category set.
\item Each \textbf{branch} is an answer; each \textbf{leaf} holds a prediction: the majority class (or class proportions)
  for classification, the mean target for regression.
\item A prediction walks from the root to one leaf. Depth $D$ means at most $D$ questions.
\end{itemize}
In feature space, each split cuts along one axis, so a tree partitions the space into axis-aligned rectangles (boxes),
each with a constant prediction.

\section{How to choose a question: impurity}
At each node we want the question that makes the two children as \emph{pure} as possible (each child mostly one class).
We need a number for impurity.

\subsection{Gini impurity}
\[ \text{Gini} = 1 - \sum_k p_k^2, \]
where $p_k$ is the fraction of class $k$ in the node. Interpretation: the probability that two items picked at random
(with replacement) from the node have different classes. Pure node: 0. Two classes 50/50: $1 - 0.25 - 0.25 = 0.5$
(the maximum for two classes).

\subsection{Entropy}
\[ H = -\sum_k p_k\log_2 p_k. \]
Interpretation: the average number of yes/no questions (bits) needed to identify the class of a random item. Pure: 0;
50/50: 1 bit; four equally likely classes: 2 bits.

\begin{example}[: impurities of one node]
8 positives, 2 negatives: $p = (0.8, 0.2)$.
Gini $= 1 - 0.64 - 0.04 = 0.32$. Entropy $= -0.8\log_2 0.8 - 0.2\log_2 0.2 = 0.8(0.322) + 0.2(2.322) = 0.258 + 0.464 = 0.722$ bits.
\end{example}

\subsection{Information gain}
The quality of a split is how much it reduces impurity, weighting each child by its share of the rows:
\[ \text{Gain} = I(\text{parent}) - \sum_{c}\frac{n_c}{n}I(c). \]

\begin{example}[: compare two splits]
Parent: 6 apply, 6 do not. Entropy $=1$.
\textbf{Split A} (``preferred city?''): left (5 yes, 1 no), right (1 yes, 5 no). Each child: $p = (5/6, 1/6)$, entropy
$= 0.650$. Weighted: $0.650$. Gain $= 1 - 0.650 = 0.350$.
\textbf{Split B} (``remote?''): left 8 rows (6 yes, 2 no), right 4 rows (0 yes, 4 no). Left entropy $= 0.811$, right $= 0$.
Weighted $= \frac{8}{12}0.811 = 0.541$. Gain $= 0.459$. \textbf{Split B wins}, even though its left child is not pure, because
the right child is perfectly pure and holds a third of the data.
With Gini: parent $0.5$; A: each child $1 - (25+1)/36 = 0.278$, gain $0.222$; B: left $1 - (0.5625 + 0.0625) = 0.375$,
weighted $0.25$, gain $0.25$. Same winner.
\end{example}

\subsection{Regression trees: variance reduction}
For a numeric target, impurity is the variance (or the sum of squared errors around the node mean). The leaf predicts the
mean.
\begin{example}[: one regression split]
Targets in a node: 2, 4, 6, 20. Mean 8; SSE $= 36 + 16 + 4 + 144 = 200$. Split ``$\{2,4,6\}$ vs $\{20\}$'': left mean 4, SSE $=
4 + 0 + 4 = 8$; right SSE $= 0$. Reduction $= 192$. The outlier gets its own leaf.
\end{example}

\subsection{How thresholds are searched}
For a numeric feature, sort the node's values; candidate thresholds are midpoints between consecutive distinct values
($m - 1$ candidates for $m$ distinct values). Evaluate the gain for every feature and every threshold; pick the best. With
running sums, all thresholds of one sorted feature can be scored in one pass. Histogram methods (LightGBM, XGBoost
\code{hist}) bucket values into about 256 bins first, making this much faster.

\section{The algorithm (pseudo-code)}
\begin{lstlisting}
def build_tree(rows, depth):
    # stopping rules
    if depth == MAX_DEPTH or len(rows) < MIN_SAMPLES_SPLIT or is_pure(rows):
        return Leaf(prediction(rows))           # majority class / mean
    best = None
    for feature in features:
        for threshold in candidate_thresholds(rows, feature):
            left, right = split(rows, feature, threshold)
            if len(left) < MIN_SAMPLES_LEAF or len(right) < MIN_SAMPLES_LEAF:
                continue
            gain = impurity(rows) - (len(left)*impurity(left) + len(right)*impurity(right)) / len(rows)
            if best is None or gain > best.gain:
                best = Split(feature, threshold, gain, left, right)
    if best is None or best.gain < MIN_GAIN:
        return Leaf(prediction(rows))
    return Node(best.feature, best.threshold,
                build_tree(best.left, depth + 1),
                build_tree(best.right, depth + 1))

def predict(node, x):
    while not node.is_leaf:
        node = node.left if x[node.feature] <= node.threshold else node.right
    return node.prediction
\end{lstlisting}
\textbf{Complexity.} Training: with pre-sorting, about $O(d\,n\log n)$ per level, so $O(d\,n\log n\cdot\text{depth})$ overall.
Prediction: $O(\text{depth})$, independent of $n$: very fast.

\subsection{Greedy, not optimal}
The algorithm picks the best split \emph{now}, never reconsidering. Finding the globally optimal tree is NP-hard. Greedy
choice can miss splits that look useless alone but are powerful together (XOR: neither $x_1$ nor $x_2$ alone reduces
impurity, but both together separate perfectly).

\section{The classic algorithms}
\begin{center}\small
\begin{tabularx}{\linewidth}{l X X X}
\toprule
& ID3 & C4.5 & CART\\
\midrule
Criterion & information gain (entropy) & gain ratio & Gini (classification), MSE (regression)\\
Splits & multiway on categoricals & multiway; handles numeric thresholds & binary only\\
Tasks & classification & classification & classification and regression\\
Missing values & no & fractional instances & surrogate splits\\
Pruning & no & error-based & cost-complexity\\
\bottomrule
\end{tabularx}
\end{center}
scikit-learn implements an optimised CART.

\subsection{The bias toward many-valued features, and gain ratio}
A feature like \code{customer\_id} with a unique value per row splits the data into pure singletons: maximum information
gain, zero generalisation. \textbf{Gain ratio} divides gain by the split information $-\sum\frac{n_c}{n}\log_2\frac{n_c}{n}$,
which is large for many small children.

\section{Overfitting and how to stop it}
A tree grown until every leaf is pure memorises the training data: training accuracy 100\%, high variance (change a few
rows and the whole structure changes, because an early split change cascades).

\textbf{Pre-pruning (early stopping)}: \code{max\_depth}, \code{min\_samples\_split}, \code{min\_samples\_leaf},
\code{max\_leaf\_nodes}, \code{min\_impurity\_decrease}, \code{max\_features}.

\textbf{Post-pruning}: grow a full tree, then cut back. \textbf{Cost-complexity pruning} (CART) minimises
\[ R_\alpha(T) = R(T) + \alpha|T|, \]
training error plus $\alpha$ times the number of leaves. Each $\alpha$ gives a nested subtree; choose $\alpha$ by cross-validation
(scikit-learn: \code{ccp\_alpha}, and \code{cost\_complexity\_pruning\_path}).
\begin{example}[: is pruning worth it?]
$\alpha = 0.02$. Subtree $T_1$: error 0.10, 5 leaves: $R = 0.10 + 0.10 = 0.20$. Pruned $T_2$: error 0.13, 2 leaves:
$R = 0.13 + 0.04 = 0.17$. Prune.
\end{example}
Post-pruning sees the whole tree, so it avoids the ``horizon effect'' of pre-pruning (stopping before a weak split that
enables strong later ones).

\section{Strengths and weaknesses}
\textbf{Strengths}: interpretable (a small tree is a flowchart a recruiter can read); handles numeric and categorical
features; no scaling needed (only the order of values matters, so any monotone transform leaves the tree unchanged);
captures interactions and non-linearity automatically; fast predictions; some implementations handle missing values.

\textbf{Weaknesses}: high variance and instability; axis-aligned boundaries approximate diagonal boundaries with
staircases; piecewise-constant predictions, so \textbf{no extrapolation} beyond the training range; greedy; biased
toward high-cardinality features; poorly calibrated probabilities from pure leaves.

\section{Feature importance}
\textbf{Impurity-based (MDI)}: sum, over all nodes splitting on feature $j$, of the weighted impurity decrease; normalise
to sum to 1. Fast, but biased toward high-cardinality and continuous features and computed on training data.
\textbf{Permutation importance}: shuffle feature $j$ in validation data and measure how much the metric drops. Model
agnostic, uses held-out data, but splits credit between correlated features. \textbf{SHAP} values give per-prediction
attributions with consistency guarantees (TreeSHAP is exact and fast for trees).

\begin{practical}[: trees at InfoEdge]
Single trees are rarely the production model, but they are (1) the unit inside every GBM ranking and lead-scoring model;
(2) a great \textbf{explanation tool}: fit a depth-3 tree to a black-box model's predictions to describe its logic to the
99acres sales team (``high intent = viewed $\ge$5 listings in one locality in 7 days and price band within 20\%'');
(3) quick \textbf{segment discovery} for the case-study round.
\end{practical}

\stuck{StatQuest: Entropy (for data science), Clearly Explained}{\ytid{YtebGVx-Fxw}}
\section{Interview questions}

\begin{qa}{\pyq{Glassdoor, InfoEdge: pseudo-code a decision tree} Write the training algorithm for a classification tree and
state its complexity.}
\oneline{Recursively, for each node, try every feature and every candidate threshold, pick the split with the largest
impurity decrease, and recurse on the two children until a stopping rule fires; training costs about
$O(d\,n\log n)$ per depth level and prediction costs $O(\text{depth})$.}
\why I would write the pseudo-code above and narrate four design decisions:
\begin{enumerate}
\item \textbf{Impurity}: Gini (default, slightly faster, no logarithm) or entropy.
\item \textbf{Candidate thresholds}: midpoints between sorted distinct values; for categorical features either one-hot
  or, for binary classification, sort categories by positive rate and treat as ordinal (provably optimal for Gini).
\item \textbf{Stopping}: max depth, min samples per split/leaf, pure node, min gain.
\item \textbf{Leaf value}: majority class and class proportions (for \code{predict\_proba}).
\end{enumerate}
Complexity: sorting each feature costs $O(n\log n)$; scanning thresholds with running class counts is $O(n)$ per feature;
at each depth level the nodes partition the $n$ rows, so total per level is $O(d\,n\log n)$ (or $O(d\,n)$ with pre-sorting
or histograms).
\pushes \emph{``How would you make it a regression tree?''} Replace Gini with variance (SSE) and the leaf value with the mean.
\emph{``How do you handle a missing value at prediction time?''} Send it down a learned default direction (XGBoost), use
surrogate splits (CART), or impute.
\end{qa}

\begin{qa}{\pyq{INT: reducing overfitting in trees} How do you stop a decision tree from overfitting?}
\oneline{Constrain growth (max depth, min samples per leaf, max leaf nodes, min impurity decrease), post-prune with
cost-complexity pruning tuned by cross-validation, or replace the single tree by an ensemble (random forest or boosting).}
\why A tree can always reach zero training error by giving every point its own leaf, which memorises noise. Each control limits
how finely the space can be cut: \code{min\_samples\_leaf=50} means no prediction is based on fewer than 50 rows, so
leaf estimates are stable. Cost-complexity pruning penalises leaves with $\alpha|T|$. Ensembles attack variance directly by
averaging (bagging) or by building in many small, shrunk steps (boosting).
\worked Shiksha lead data: unlimited depth 100\% train / 71\% val; \code{max\_depth=8, min\_samples\_leaf=40}: 83\% / 80\%;
random forest: 88\% / 85\%.
\pushes \emph{``Which is better, pre- or post-pruning?''} Post-pruning is usually better because it sees later useful splits;
pre-pruning is cheaper. In practice, tune \code{min\_samples\_leaf} and \code{max\_depth} by CV.
\end{qa}

\begin{qa}{\pyq{OA: entropy and information gain} A node has 8 positives and 2 negatives. Compute its entropy and Gini. Then
compute the information gain of splitting a 6/6 parent into (5,1) and (1,5).}
\oneline{Entropy $0.722$ bits, Gini $0.32$; information gain $1 - 0.650 = 0.350$ bits.}
\why As in the worked examples: $H(5/6,1/6) = -\frac56\log_2\frac56 - \frac16\log_2\frac16 = 0.219 + 0.431 = 0.650$; both
children have equal size and entropy, so the weighted child entropy is 0.650.
\end{qa}

\begin{qa}{Gini or entropy: does it matter?}
\oneline{Rarely: they choose the same split in the vast majority of cases; Gini is slightly cheaper (no logarithm), entropy
is slightly more sensitive to changes in small class probabilities.}
\why Both are concave, zero for pure nodes and maximal for uniform ones; Gini is roughly a scaled second-order approximation
of entropy. Empirical studies find them disagreeing on about 2\% of splits. Mention \textbf{misclassification error}
$1 - \max_kp_k$ as a third option that is a poor splitting criterion: it is not strictly concave, so it often sees zero gain
for useful splits.
\end{qa}

\begin{qa}{\deep A regression tree is trained on flats priced 20 to 90 lakh. What does it predict for a mansion that is far
larger than anything in training? Can a random forest or gradient boosting do better?}
\oneline{It predicts the mean of whichever leaf the mansion falls into, at most about 90 lakh: trees cannot extrapolate beyond
the range of training targets, and forests and GBMs (sums or averages of such leaves) cannot either.}
\why Each leaf predicts a training-data average, and every split beyond the largest training value sends the mansion to the
rightmost leaf. Random forest predictions are averages of leaf means (bounded by the target range). Gradient boosting
sums many trees' leaf values; in practice it also saturates. For trend extrapolation (prices rising over time, sales growing),
detrend first and model residuals with the tree, use linear models or linear trees (LightGBM \code{linear\_tree}), or add a
linear component.
\worked Training max price 90 lakh. A 5{,}000 sq ft villa: tree says 88 lakh (the mean of the largest-flat leaf). A linear model
on log price extrapolates to about 3 crore.
\pushes \emph{``Where does this bite in practice?''} Forecasting with time trends: a GBM trained on 2023--2025 application
counts cannot predict 2026 values above the historic maximum. Use differencing or a ratio target.
\end{qa}

\begin{qa}{\deep Why do decision trees not need feature scaling? Which preprocessing \emph{can} change a tree?}
\oneline{Splits depend only on the order of values, so any strictly increasing transform (scaling, log) leaves the tree unchanged;
changes that alter order or grouping (one-hot vs ordinal encoding of categories, binning, combining features) do change it.}
\why A threshold $x\le t$ on $x$ is identical to $\log x \le \log t$. But ordinal-encoding ``city'' as 1, 2, 3, \dots\ imposes an
arbitrary order, so the tree can only split cities into contiguous blocks of that order. One-hot gives one city versus the rest per
split. Creating a ratio feature (price per sq ft) can let one split do what needed many.
\end{qa}

\begin{qa}{Why do trees prefer features with many unique values, and how is it fixed?}
\oneline{More unique values offer more candidate splits, so by chance alone one of them looks good, and an ID-like feature can
produce pure singleton leaves; gain ratio, minimum leaf sizes, permutation importance and dropping IDs fix it.}
\why With $m$ distinct values there are $m-1$ thresholds to try; the maximum of many noisy gains is optimistic (multiple comparisons).
ID3 with \code{customer\_id} gives maximum information gain and zero generalisation. Fixes: C4.5 gain ratio, conditional inference
trees (statistical tests), \code{min\_samples\_leaf}, and never feeding IDs.
\end{qa}

\begin{qa}{\deep Can a decision tree learn XOR? What does this tell you about greedy splitting?}
\oneline{A depth-2 tree can represent XOR, but a greedy learner may not find it, because no single first split reduces impurity
on its own; with real data's noise it usually finds it, but in pure form the first split's gain is zero.}
\why Four points $(0,0)\to0$, $(0,1)\to1$, $(1,0)\to1$, $(1,1)\to0$: splitting on $x_1$ gives children (0,1) and (1,0), each
50/50: zero gain. Same for $x_2$. A greedy algorithm with a ``min gain $> 0$'' rule stops at the root. Without that rule it splits
anyway (ties broken arbitrarily) and the second level separates perfectly. Lesson: greedy, one-step look-ahead misses
interactions whose components are individually useless; ensembles with random feature choices and deeper trees mitigate it.
\end{qa}

\begin{qa}{How does a decision tree handle categorical features?}
\oneline{CART-style implementations need numeric input, so categories are one-hot or ordinal encoded; some libraries split
categories natively by grouping them into two subsets, using the trick of sorting categories by target mean.}
\why For binary classification or regression, sorting the $k$ categories by mean target and treating them as ordered finds the best
two-way partition among $2^{k-1}-1$ possibilities in $O(k\log k)$. LightGBM and CatBoost do native categorical splits; recent
scikit-learn also supports them in \code{HistGradientBoosting}. One-hot with many levels makes deep, unbalanced trees.
\end{qa}

\begin{qa}{How does a tree compute \code{predict\_proba}, and why are those probabilities poorly calibrated?}
\oneline{It returns the class proportions of the training rows in the leaf; deep trees have small, pure leaves that output 0 or 1,
so probabilities are overconfident.}
\why A leaf with 3 positives out of 3 says $p = 1.0$, but the true rate for such inputs might be 0.7. Remedies: larger
\code{min\_samples\_leaf}, Laplace smoothing $(k+1)/(n+2)$, ensembling (averages are smoother), and post-hoc calibration
(isotonic/Platt).
\end{qa}

\begin{qa}{What is cost-complexity pruning? Work an example.}
\oneline{It chooses the subtree that minimises training error plus $\alpha$ times the number of leaves, with $\alpha$ tuned by
cross-validation.}
\why Increasing $\alpha$ from 0 gives a sequence of nested subtrees (weakest-link pruning): at each stage prune the internal node
whose removal increases error least per leaf removed. Worked example in the theory section: $0.20$ vs $0.17$ means prune.
\end{qa}

\begin{qa}{Explain impurity-based vs permutation feature importance. Which would you trust?}
\oneline{Impurity importance sums each feature's split gains on training data and is biased toward high-cardinality features;
permutation importance measures the validation metric drop when a feature is shuffled and is usually more trustworthy, though
it splits credit among correlated features.}
\why A random ID column gets high MDI importance (many splits that fit noise) but near-zero permutation importance (shuffling it
does not hurt validation). Correlated features: shuffling one leaves its twin intact, so both look unimportant. Fix by grouping
correlated features or using SHAP with care.
\end{qa}

\begin{qa}{\deep Why are single decision trees called unstable, and why is that a good thing for ensembles?}
\oneline{A small change in the data can change an early split and therefore the whole tree below it; that high variance means
different trees make different errors, which averaging cancels.}
\why Bagging's variance reduction is $\rho\sigma^2 + \frac{1-\rho}{B}\sigma^2$: it needs models that individually have high
variance $\sigma^2$ and are not too correlated ($\rho$ small). Deep trees have low bias and high variance: ideal for bagging.
Stable models like linear regression gain little from bagging.
\end{qa}

\begin{qa}{What does a decision boundary of a tree look like in 2D? What kind of boundary is hard for it?}
\oneline{A set of axis-aligned rectangles; diagonal or curved boundaries need many small steps (a staircase), requiring a deep tree.}
\why A boundary like $x_1 = x_2$ is approximated by a staircase; with limited depth the approximation is coarse. Fixes: add a
feature $x_1 - x_2$ (one split then suffices), use oblique trees, or use ensembles (many staircases averaged look smoother).
\end{qa}

\begin{qa}{How many leaves can a binary tree of depth 10 have? How many candidate thresholds for a feature with 100
distinct values?}
\oneline{At most $2^{10} = 1024$ leaves; 99 candidate thresholds.}
\end{qa}

\begin{qa}{\deep You fit a depth-3 tree to explain a complex model's predictions to a business team. What is this called and what
are its risks?}
\oneline{A global surrogate model; the risk is that it explains the surrogate, not the original, so you must report its fidelity
(how often it agrees with the black box) and keep it shallow enough to read.}
\why Fit the tree to the black-box predictions (not the true labels). If fidelity is 92\%, say so; regions of disagreement deserve
local explanations (SHAP). Useful for case-study rounds: ``this is the logic the lead scorer learned''.
\end{qa}

\begin{qa}{Compute the Gini impurity of a node with class counts (2, 3, 5) and the entropy of one with counts (3, 3, 6).}
\oneline{Gini $= 1 - (0.04 + 0.09 + 0.25) = 0.62$; entropy $= 1.5$ bits.}
\why Entropy: $p = (0.25, 0.25, 0.5)$; $-2(0.25\log_20.25) - 0.5\log_20.5 = 2(0.5) + 0.5 = 1.5$.
\end{qa}

\begin{qa}{How do trees handle missing values?}
\oneline{Options: impute first; CART surrogate splits (use a correlated feature when the split feature is missing); C4.5 sends the
row down all branches with fractional weights; XGBoost/LightGBM learn a default direction for missing values at each split.}
\why XGBoost tries sending all missing rows left, then right, and keeps the better gain; this also captures ``missingness is
informative'' (salary blank often means fresher). scikit-learn's \code{DecisionTreeClassifier} supports NaN in recent versions;
\code{HistGradientBoosting} always has.
\end{qa}

\begin{qa}{\deep A tree's top split is on ``days since last login''. The product manager concludes recent login causes applying.
How do you respond?}
\oneline{A split shows predictive association, not causation; recent login may simply mark already-active job seekers, so I would
propose an experiment (re-engagement nudges with a holdout) to measure the causal effect.}
\why Trees pick the most predictive feature, and the most predictive feature is often a proxy for intent. Causal questions need
interventions or causal designs. This is a common case-study trap.
\end{qa}
